\clearpage\begingroup\SourceLeaf{651b}
\ArithmeticHead{16}{L}{H. GEPPERT, ÜBER GAUSS’ ARBEITEN ZUR MECHANIK UND POTENTIALTHEORIE.}
\makeatletter\tagsleft@true\let\veqno\@@leqno\makeatother
\renewcommand{\thefootnote}{\arabic{footnote})}\setcounter{footnote}{28}
\begin{center}{\bfseries II. Das Prinzip des kleinsten Zwanges.}\\[5pt]
\textls[100]{1. Historisches.}\end{center}
\noindent\hspace*{\parindent}Das Prinzip des kleinsten Zwanges ist ebenso wie das Prinzip der kleinsten
Quadrate bei \textsc{Gauss} in seiner philosophisch-ethischen Weltanschauung begründet.
\textsc{Sartorius von Waltershausen}\footnote{W. \textsc{Sartorius von Waltershausen}, \textit{Gauss zum Gedächtnis.} Leipzig 1856. S. 97.} schildert \textsc{Gauss} als eine in ihrem
geistigen, wie in ihrem Gemütsleben durchaus wahre Natur und fährt fort:
»Dieser Durst nach Wahrheit verbunden mit einem heiligen Drang nach Gerechtigkeit
bezeichneten vornehmlich \textsc{Gauss}’ erhabenen Charakter. Beide Gefühle
schlummerten von jeher in der Tiefe seines Wesens, sie hingen auf das
Innigste mit seinen philosophischen und mit seinen religiösen Betrachtungen
zusammen und sind ohne Frage durch seine erhabene Naturbetrachtung noch
weiter ausgebildet und bekräftigt worden. Das Prinzip des kleinsten Zwanges
war gleichsam die mathematische Verkörperung jenes ethischen Grundgedankens,
\clearpage\SourceLeaf{652}\ArithmeticHead{17}{R}{II. DAS PRINZIP DES KLEINSTEN ZWANGES.}
\HistoricalFoot{X 2 Abh. 7.}{3}
\noindent den er für das Universum als bindend erkannte.« Mit dieser Darstellung
stimmt es überein, wenn \textsc{Gauss} selbst am Schlusse seiner Abhandlung\footnote{\textit{Über ein neues allgemeines Grundgesetz der Mechanik.} Journal für reine u. angew. Mathematik,
4, 1829, S. 232; Werke V, 1867, S. 23—28. Übersetzt und herausgeg. in \textsc{Ostwalds} Klassiker, Bd. 167,
1908, von \textsc{Philip E. B. Jourdain}.}, in
der er sein Prinzip formulierte, es als merkwürdig bezeichnet, dass die Natur
die freien Bewegungen eines mechanischen Systems unter dem Druck von
Bedingungen ebenso modifiziert, wie der rechnende Mathematiker Erfahrungen
über untereinander abhängige Grössen ausgleicht; denn letzterer verfährt ja
eben nach einem »Prinzip der Gerechtigkeit«.

Neben diese ethische Auffassung des \textsc{Gauss}schen Prinzips tritt die erkenntnistheoretische.
Für \textsc{Gauss} handelt es sich in erster Linie um ein Ökonomieprinzip
im \textsc{Mach}schen Sinne: er selbst nennt als den Hauptvorzug des
neuen Prinzips, dass es in gleicher Weise Statik und Dynamik in denkbar
grösster Allgemeinheit umfasse, die Trennung zwischen dem \textsc{d’Alembert}schen
Prinzip und dem der virtuellen Verrückungen überflüssig mache, und von
höherem Standpunkte aus die Quintessenz aller mechanischen Erfahrungen
bilde. Dabei hat \textsc{Gauss} zu jener Zeit die Tragweite seines Prinzips geringer
eingeschätzt, als sie wirklich ist; nach unserer heutigen Auffassung ist es das
umfassendste aller mechanischen Prinzipien. \textsc{Gauss} selbst hat, wie er ausdrücklich
sagt, zunächst im Auge, einen einleuchtenderen, das Creditiv gewissermassen
in sich selbst tragenden Grundsatz zu gestalten, der die beiden
oben genannten Prinzipien in sich schliesst, ohne ihnen von vornherein etwas
Neues hinzufügen zu wollen. Die Form seines Gesetzes erklärt sich einerseits
daraus, dass \textsc{Gauss} sich schon seit 1794 intensiv mit der Methode der kleinsten
Quadrate beschäftigte und wohl eine methodisch analoge Form der Naturgesetze
suchte — wie ja denn auch diese Analogie in der \textsc{Gauss}schen Formulierung
eine entscheidende Rolle spielt\footnote{Vgl. Brief an \textsc{Olbers} vom 31. I. 1829, Werke XI, 1, 1927, S. 22, und die von \textsc{Gauss} beeinflusste
Dissertation: A. \textsc{Ritter}, \textit{Über das Prinzip des kleinsten Zwanges.} Göttingen 1853, S. 23 ff.}, — andererseits daraus, dass die
Tatsache eines Minimumprinzips der an sich zwar überwundenen, aber in den
Grundlagen der Mechanik noch nachschwingenden teleologischen Naturauffassung
Rechnung trägt.

Bei \textsc{Leibniz} und \textsc{Euler} ist die Teleologie geradezu bestimmend, und man
ist im 16. und 17. Jahrhundert immer geneigt, in den physikalischen Ge\nobreakdash-
\clearpage\SourceLeaf{653}\ArithmeticHead{18}{L}{H. GEPPERT, ÜBER GAUSS’ ARBEITEN ZUR MECHANIK UND POTENTIALTHEORIE.}
\noindent setzen besonders weise Anordnungen eines denkenden Schöpfers zu erblicken;
\textsc{Euler} spricht es in seiner Methodus inveniendi von 1741 ganz offen aus, man
könne die physikalischen Phänomene nicht nur aus ihren Ursachen, sondern
auch aus ihrem Zweck erklären, indem man nichts in der Welt anträfe, woraus
nicht irgendeine Maximums- oder Minimumseigenschaft hervorleuchte\footnote{Vergl. E. \textsc{Mach}, \textit{Die Mechanik in ihrer Entwicklung.} Leipzig 1883, S. 428. E. \textsc{Dühring},
\textit{Kritische Geschichte der allgemeinen Prinzipien der Mechanik.} Berlin 1873, S. 292—296.}.
Ihren Höhepunkt erreichte die teleologische Auffassung in dem von \textsc{Maupertuis}
ausgesprochenen Prinzip der kleinsten Aktion, das wegen seiner unsauberen Formulierung
unklaren metaphysischen Spekulationen Tür und Tor öffnete. \textsc{Lagrange}
war es, der diesen unwissenschaftlichen Überlegungen ein Ende setzte; er sieht
den Wert der Prinzipien nach seiner, unserer heutigen nahestehenden Auffassung
in der Ökonomie des Denkens und der Kondensation aller mechanischen Erfahrung
in eine Regel, wobei er ausdrücklich erklärt, er wolle von allen, doch sehr prekären
theologischen und metaphysischen Überlegungen absehen. Gleichzeitig entkleidete
er das Prinzip der kleinsten Aktion seines metaphysischen Beigeschmacks,
indem er es in eine exakte Form brachte und unter der Annahme skleronomer
Bedingungsgleichungen aus den anderen Grundgesetzen bewies. Bei \textsc{Lagrange}
selbst erhielt dadurch das \textsc{Maupertuis}sche Prinzip eine untergeordnete Rolle, und
die teleologische Spekulation war damit für die Folgezeit erledigt. Aber die
geschichtliche Entwicklung der Mechanik im 18. Jahrhundert spiegelt sich in
der Folgezeit, wie noch heute in der Behandlung der Elemente, so damals in der
Fassung ihrer Hauptsätze wieder. Und so ist es auch für \textsc{Gauss} ein Gewinn an
Einsicht, wenn er seinem Prinzip den Stempel eines Minimumsatzes aufprägt:
»Die Natur verhält sich so, dass sie alle Bewegungen unter möglichst kleinem
Zwange ablaufen lässt.« In Wahrheit kann es sich bei dieser Form garnicht
um ein teleologisches Prinzip im eigentlichen Sinne handeln, da die wirklich
stattfindende Bewegung des mechanischen Systems nur gegenüber fingierten,
eigentlich unmöglichen Bewegungen, aber nicht gegenüber anderen ebenfalls
möglichen Zuständen ausgezeichnet wird\footnote{Vergl. A. \textsc{Voss}, \textit{Die Prinzipien der rationellen Mechanik.} Encykl. d. math. Wiss. IV, 1, S. 19 ff.}. Schliesslich ist beim \textsc{Gauss}schen
Prinzip zu beachten, dass es sich nicht, wie bei dem der kleinsten Aktion, um
ein Variationsprinzip handelt, sondern nur um eine Minimumseigenschaft der
jeweils augenblicklichen Beschleunigungen.
\clearpage\SourceLeaf{654}\ArithmeticHead{19}{R}{II. DAS PRINZIP DES KLEINSTEN ZWANGES.}
\HistoricalFoot{}{3\textsuperscript{*}}
Ist so die Auffindung des Prinzips vom kleinsten Zwange für die Beurteilung
von \textsc{Gauss}’ Persönlichkeit in hohem Grade charakteristisch, so ist
sie andererseits in der Arbeit seiner grossen Vorgänger \textsc{Euler}, \textsc{d’Alembert},
\textsc{Lagrange} verwurzelt. Nach der damaligen allgemeinen Auffassung spaltet sich
die Mechanik in die beiden, noch völlig getrennten Disziplinen der Statik
und Dynamik auf, deren erste durch das Prinzip der virtuellen Verrückungen,
deren zweite durch das \textsc{d’Alembert}sche Prinzip beherrscht wird.

Das Prinzip der virtuellen Verrückungen war in Spezialfällen schon von
\textsc{Stevin}, \textsc{Galilei} und \textsc{Torricelli} bemerkt worden, die die Beobachtung machten,
dass das Gleichgewicht bei den einfachen Maschinen auf allgemeinen Sätzen
beruhe; aber erst \textsc{Joh. Bernoulli} gestaltet hieraus in genialer Induktion das
eigentliche Prinzip der virtuellen Geschwindigkeiten\footnote{In einem Briefe an \textsc{Varignon} vom 26. Januar 1717, vergl. P. \textsc{Varignon}, \textit{Nouvelle Mécanique}
2 vol. Paris 1725. Wegen der ausführlichen Zitate hier und im folgenden vergl. \textsc{Voss} a.\,a.\,O., \textsc{Mach} a.\,a.\,O.}. \textsc{Lagrange} machte in
der 1. Auflage seines Hauptwerkes (1788)\footnote{J. L. \textsc{Lagrange}, \textit{Mécanique analytique.} Paris 1788; Oeuvres complètes, vol. 11, p. 23.} aus diesem Satz das Grundprinzip
der Statik, und zwar auf dem Wege axiomatischer Einführung, die er später
in der 2. Auflage (1811) durch einen Beweisversuch ersetzte, der die an einfachen
Maschinen (Flaschenzügen) gewonnenen Anschauungen für die allgemeinen
mechanischen Systeme nutzbar machte.

Das genannte Prinzip können wir folgendermassen fassen. Es sei $\Sigma$ ein
mechanisches System von $n$ Massenpunkten, deren kartesische Koordinaten
wir uns durchgezählt denken und im folgenden mit $x_i\ (i=1,\, .\, .\, .\, ,\,3n)$ bezeichnen,
sodass $x_1,x_2,x_3$ die Koordinaten des ersten, $ .\, .\, .\, x_{3n-2},x_{3n-1},x_{3n}$
die des $n$-ten Massenpunktes bedeuten. Auf die $n$ Punkte mögen ebensoviele
Kräfte wirken, deren Komponenten in entsprechender Weise mit $X_i\ (i=1,\, .\, .\, .\, 3n)$
bezeichnet seien. Ferner möge das System gewissen $m$ Bedingungsgleichungen
unterworfen sein, die wir in allgemeinster differentieller, etwa rheonomer
Form schreiben können:
\[
\sum_{i=1}^{3n}(a_{\varkappa i}\,dx_i+a_{\varkappa}\,dt)=0,
\qquad(\varkappa=1\, .\, .\, .\, m).\tag{1}
\]
\noindent Das Prinzip der virtuellen Verrückungen besagt dann: Notwendig und hinreichend
dafür, dass $\Sigma$ sich im Gleichgewicht befindet, ist das Bestehen der
\clearpage\SourceLeaf{655}\ArithmeticHead{20}{L}{H. GEPPERT, ÜBER GAUSS’ ARBEITEN ZUR MECHANIK UND POTENTIALTHEORIE.}
\noindent Gleichung
\[
\sum_{i=1}^{3n}X_i\delta x_i=0,\tag{2}
\]
\noindent worin $\delta x_i$ die Komponenten von virtuellen, d.\,h. durch die Bedingungsgleichungen
gestatteten infinitesimalen Verschiebungen der Massenpunkte sind,
nämlich von solchen, die den Gleichungen
\[
\sum_{i=1}^{3n}a_{\varkappa i}\delta x_i=0,
\qquad(\varkappa=1\, .\, .\, .\, m)\tag{3}
\]
\noindent genügen\footnote{Diese allgemeinste Fassung des Prinzips s.\,z.\,B. bei \textsc{Cl. Schaefer}, \textit{Die Prinzipe der Dynamik.}
Berlin 1919.}. Bei \textsc{Lagrange} sind nur Bedingungen berücksichtigt, die durch
Gleichungen zwischen den Punktkoordinaten realisiert werden; die dann möglichen
Verschiebungen $\delta x_i$ sind — und ebenso bei anholonomen Bedingungen
— von der Art, dass zu jeder auch die entgegengesetzte denkbar ist. Als
allgemeines Gleichgewichtskriterium schliesst das genannte Prinzip die ganze
Statik in sich.

Den tragenden Grundpfeiler der Dynamik bildet hingegen das Prinzip
von \textsc{d’Alembert}\footnote{J. \textsc{d’Alembert}, \textit{Traité de Dynamique.} Paris 1743.}. Bezeichnen wir die Massen der betrachteten Punkte analog
wie die Koordinaten mit $m_i\ (i=1,\, .\, .\, .\, 3n)$, sodass man $m_1=m_2=m_3,\, .\, .\, .\,
m_{3n-2}=m_{3n-1}=m_{3n}$ zu setzen hat, und sind die $X_i$ wieder die Komponenten
der auf die Massen $m_i$ wirkenden Kräfte, so erleiden die Punkte unter dem Einfluss
der genannten Kräfte einerseits und der von den Bedingungsgleichungen
herrührenden Zwangskräfte andererseits die Beschleunigungen $\ddot x_i$; dann sind die
sog. verlorenen Kräfte $X_i-m_i\ddot x_i$ mit den durch die Bindungen hervorgerufenen
Zwangskräften im Gleichgewicht. Dies ist der \textsc{d’Alembert}sche Satz, den man
bis zu einem gewissen Grade als rein logische Überlegung ansprechen kann.
\textsc{Lagrange} tat nun in seiner Mécanique den naheliegenden Schritt, diesen Satz
mit dem Prinzip der virtuellen Verrückungen zu verbinden, sodass bei ihm
das \textsc{d’Alembert}sche Prinzip die Gestalt erhält
\[
\sum_{i=1}^{3n}(X_i-m_i\ddot x_i)\delta x_i=0,\tag{4}
\]
\clearpage\SourceLeaf{656}\ArithmeticHead{21}{R}{II. DAS PRINZIP DES KLEINSTEN ZWANGES.}
\noindent worin die $\delta x_i$ ein System virtueller Verrückungen im obigen Sinne bezeichnen,
die den Gleichungen (3) genügen. Der \textsc{d’Alembert}sche Satz führt somit die
Dynamik auf die Statik zurück und enthält letztere als Spezialfall, in dem
$\ddot x_i=0$ zu setzen ist. Er ist äquivalent mit den \textsc{Lagrange}schen Grundgleichungen
der Mechanik.
\clearpage\endgroup
